% Integrated candidate; input from cas-sc-template_F05_applied.tex.
% Archived measurements remain descriptive; audit restrictions are explicit.
\section{Controlled Experiments}
\label{sec:controlled_experiments}

\subsection{Implementation identity and experimental design}
The archived experiments comprise 340 two-flow runs: 60 each for the
Ubuntu in-tree baseline, the separately loaded BBRv1 module, BBR2, and
BBRv3; 80 E3 runs; and 20 artificial-loss runs. The manifests identify
BBRv3 by the 6.13.7-bbr3 kernel, source marker BBR\_VERSION=3, and runtime
name \texttt{bbr}. The BBRv1 module registers as \texttt{bbr1} and uses
a recorded compatibility patch for an unexported TSO helper. BBR2 uses
a separate 5.13.12-bbr2alpha kernel and runtime name \texttt{bbr2}.
The in-tree baseline remains a distinct implementation label.
Different kernel environments and the module patch limit causal
attribution of performance differences to algorithm version alone.

The E1/E2 matrix uses a 100-Mbps bottleneck, a 20-ms first-flow RTT,
second-flow RTTs of 20, 40, or 80 ms for the same-algorithm comparison,
and equal-RTT BBR/CUBIC competition. Queue targets are 0.5, 2, or 10
times the reference BDP. Each condition has five 180-s repetitions,
with ECN disabled. Receiver-side iperf3 summaries provide the archived
full-run goodput. Jain fairness and flow shares are derived from those
per-flow rates.

The configuration declares a 30-s warmup, but the primary archived
analysis does not exclude it. A separate sensitivity analysis uses
complete receiver intervals within the configured 30--180-s window,
without estimating bytes in boundary intervals. Across 68 condition
summaries, the largest absolute changes in mean total goodput, Jain
index, and second-flow share are 1.313 Mbps, 0.02949, and 0.01678,
respectively. These are sensitivity results rather than revised
acceptance thresholds.

Loaded RTT measurements are sequential early ping samples from the two
senders. Queue counters were sampled before the transfers completed.
They therefore do not share the goodput observation window, and a zero
counter in a saved snapshot does not establish zero drops over 180 s.
The experiment records cannot reconstruct the missing end-of-run
queue counters or a simultaneous post-warmup RTT window.

% Candidate fragment generated from 06_执行台账/cross_version_20261012/.
% Verify figure filenames and section numbering against the Overleaf main version.
\subsection{Conditioned cross-version comparison}
The baseline, an independently loaded BBRv1 module, BBR2, and BBRv3 runs were analyzed with the same condition
keys, full-run receiver-summary window, and five-run repetition rule.
Although configuration files declare a 30-s warmup, the archived primary
analysis uses the full 180-s transfer. A separate interval-based
sensitivity analysis retains the requested post-warmup window. The Ubuntu in-tree
baseline is retained as an implementation-specific reference. BBRv1,
BBR2, and BBRv3 are identified by their source/kernel markers; their runtime
registration names are \texttt{bbr1}, \texttt{bbr2}, and \texttt{bbr},
respectively. The BBRv1 module uses the recorded 6.13 compatibility patch
for the unexported TSO helper.

Across the tested 100-Mbps matrix, equal-RTT total goodput remains close to
the bottleneck for all four implementations, while flow shares and Jain
fairness vary with RTT asymmetry, queue setting, and implementation. Under
the 40- and 80-ms second-flow conditions, the observed long-RTT share differs across implementations. These results support a
conditional comparison over the tested matrix; they do not establish a
universal version ranking or identify unmeasured internal state.

For the no-ECN BBR/CUBIC comparison, total goodput is similar across the
four implementation environments, while the BBR/CUBIC share and Jain
index remain condition dependent. The corresponding values and raw inputs
are archived in \texttt{cross\_version\_comparison.csv}; figures are generated by
\texttt{analyze\_cross\_version.py}.

% Candidate fragment generated from E3 artifacts on 2026-10-12.
% Do not merge into Overleaf until authors verify section numbering, citations, and figure format.
% Evidence: 06_执行台账/formal_e3_20261011/condition_summary.csv
\subsection{Additional comparison under equal RTT}
The additional experiment matrix contains eight network cases at 50 or
200~Mbps, 10 or 60~ms RTT, and two frozen buffer labels.
For the 50-Mbps cases, the configured absolute buffers correspond
to 4 and 16 times the physical reference BDP, although the frozen file
labels them 1 and 4. For the 200-Mbps cases, those labels are consistent
with the capacity-times-reference-RTT BDP. The original predictions
are retained; the affected cases do not validate the intended buffer bins. We ran two BBRv3 flows
and two CUBIC flows separately for five independent repetitions per case
(80 runs in total). ECN is disabled in all cases.

Table~\ref{tab:e3_candidate} reports mean total goodput and Jain fairness.
The values are condition-specific observations rather than universal
version rankings. The BBRv3 fairness values were close to one in these equal
RTT cases, while CUBIC showed more variation. The queue snapshots do not cover the complete measurement window;
this experiment therefore does not establish full-run congestion-loss
or retransmission behavior.
The archived goodput and Jain values use full-run receiver summaries.
A separate post-warmup interval sensitivity analysis is retained;
the table does not report warmup-excluded steady-state estimates.

\begin{table}[htbp]
\centering
\caption{Descriptive E3 full-run means. Goodput (Mbps) and Jain cells report BBRv3/CUBIC. C is Mbps and RTT is ms. Buffer uses decimal kB; actual target Q/BDP is computed from capacity and reference RTT, before packet-limit rounding. Frozen labels are preserved in the input audit.}
\label{tab:e3_candidate}
\begin{tabular}{lrrrrrr}
\toprule
Case & C & RTT & Buffer (kB) & Actual Q/BDP & Goodput & Jain \\
\midrule
E3\_001 & 50 & 10 & 250 & 4 & 47.1/47.8 & 1.000/0.969 \\
E3\_002 & 50 & 10 & 1000 & 16 & 47.2/47.8 & 1.000/0.917 \\
E3\_003 & 50 & 60 & 1500 & 4 & 46.9/47.7 & 1.000/0.928 \\
E3\_004 & 50 & 60 & 6000 & 16 & 46.8/47.7 & 1.000/0.987 \\
E3\_005 & 200 & 10 & 250 & 1 & 188.2/191.2 & 0.995/0.963 \\
E3\_006 & 200 & 10 & 1000 & 4 & 188.0/191.2 & 0.999/0.955 \\
E3\_007 & 200 & 60 & 1500 & 1 & 187.0/190.8 & 0.997/0.963 \\
E3\_008 & 200 & 60 & 6000 & 4 & 187.1/190.8 & 0.996/0.938 \\
\bottomrule
\end{tabular}
\end{table}


\subsection{Artificial loss and experimental limitations}
The artificial-loss extension compares BBRv3 and CUBIC at configured
loss probabilities of zero and 0.1\% on a separate access link,
using five runs per algorithm and setting. Archived full-run mean total
goodput changes from 93.87 to 94.25 Mbps for BBRv3 and from 95.58 to
81.06 Mbps for CUBIC. The small BBRv3 increase is a descriptive observation
and is not evidence that loss improves performance. The loss generator's
seed was not exposed in the retained configuration; this setup does not
represent correlated wireless or satellite errors.

The additional equal-RTT cases extend the capacity and RTT settings of
the main matrix. Their physical queue sizes are reported explicitly in
Table~\ref{tab:e3_candidate}. Because four cases differ from the intended
buffer labels, we use these results as an additional empirical comparison,
rather than as a validation of the original buffer-category predictions.
The experiments evaluate goodput and bandwidth sharing under controlled
conditions; they do not establish universal suitability thresholds or
isolate the causal effect of an individual controller mechanism.
